\section{Ścieżki eksperymentalne}

W ramach Hearthstone AI Competition silnik obsługuje dwie niezależne ścieżki eksperymentalne, stawiając przed agentami odmienne wyzwania \cite{hearthstone_competition}.

Pierwsza ścieżka, czyli kategoria przygotowanych talii (ang. \textit{premade deck playing track}), polega na graniu z góry określonymi taliami, których składy są także wcześniej znane --- agent zna listę kart swojej talii i talii przeciwnika. Nie oznacza to jednak pełnej obserwowalności przebiegu partii, ponieważ podczas gry agent nadal nie zna bieżącej ręki przeciwnika ani kolejności kart pozostałych w obu taliach. Zadaniem algorytmu jest wypracowanie jak najlepszej strategii rozgrywania akcji przy ustalonej puli kart.

Druga ścieżka, czyli kategoria talii stworzonych przez użytkowników (ang. \textit{user created deck playing track}), przyznaje agentowi pełną swobodę w doborze talii. Talia może być zarówno przygotowana ręcznie przez uczestników konkursu, jak i wypracowana automatycznie przez dowolny algorytm. Optymalizacja składu talii staje się tu integralną częścią problemu i jest równie istotna jak strategia samej rozgrywki, co czyni tę ścieżkę szczególnie interesującą.
